\chapter{Your laptop}
\label{ch:laptop}

\chapterbrief
  {Install Java, Maven and git, and prove each one works.}
  {Package managers, and why version numbers matter here.}
  {Three commands that print version numbers.}

Everything in this chapter happens on \textbf{your own computer}, not the Pi.

\section{A package manager first}

Installing software by downloading installers from websites is how you end up
with three Javas and no idea which one runs. Use a package manager: one command
to install, one to update, one to remove.

\subsection*{macOS --- Homebrew}

\begin{shell}
# on your laptop
/bin/bash -c "$(curl -fsSL \
  https://raw.githubusercontent.com/Homebrew/install/HEAD/install.sh)"
\end{shell}

It will ask for your password and print a couple of commands to run afterwards
(adding \code{brew} to your \code{PATH}) --- \textbf{actually run them}, then
close and reopen your terminal.

\subsection*{Linux (Debian/Ubuntu)}

Already installed: \code{apt}.

\subsection*{Windows}

Two options, and this one genuinely matters:

\begin{trap}{Use WSL2, not native Windows}
Install \textbf{WSL2 with Ubuntu} (\code{wsl --install} in an Administrator
PowerShell, then reboot). Do all your work inside that Ubuntu.

Native Windows will fight you on line endings, file permissions, SSH paths and
shell scripts --- four separate categories of confusion that have nothing to do
with what you are trying to learn. Inside WSL2, every command in this book works
exactly as printed.

Keep your project files under the Linux home directory (\code{\textasciitilde/})
and \emph{not} under \code{/mnt/c/}, or filesystem performance will be dreadful.
\end{trap}

\section{Java 21}

\begin{shell}
# macOS
brew install openjdk@21

# Debian / Ubuntu / WSL2
sudo apt update && sudo apt install -y openjdk-21-jdk
\end{shell}

On macOS, Homebrew does not link the JDK where the system expects it. Run the
command it prints, which looks like:

\begin{shell}
sudo ln -sfn /opt/homebrew/opt/openjdk@21/libexec/openjdk.jdk \
             /Library/Java/JavaVirtualMachines/openjdk-21.jdk
\end{shell}

Verify:

\begin{shell}
java -version
\end{shell}

\begin{trace}
openjdk version "21.0.11" 2026-04-21
OpenJDK Runtime Environment (build 21.0.11+10)
OpenJDK 64-Bit Server VM (build 21.0.11+10, mixed mode)
\end{trace}

\begin{concept}{Why 21 specifically}
Java 21 is a \emph{Long Term Support} release, and it has one feature this
project uses constantly: \code{HexFormat}, for parsing and printing hexadecimal
without pulling in a library.

Java 17 would mostly work. Java 11 will not compile the code in Volume One.
Newer than 21 is fine.

If \code{java -version} shows something older, you have two Javas installed and
the wrong one is first in your \code{PATH} --- see
Appendix~\ref{app:trouble0}.
\end{concept}

\section{Maven}

\begin{shell}
# macOS
brew install maven

# Debian / Ubuntu / WSL2
sudo apt install -y maven
\end{shell}

\begin{shell}
mvn -v
\end{shell}

\begin{trace}
Apache Maven 3.9.6
Maven home: /opt/homebrew/Cellar/maven/3.9.6/libexec
Java version: 21.0.11, vendor: Homebrew
\end{trace}

Check that ``Java version'' line --- Maven has its own idea of which Java to use,
and if it says 17 while \code{java -version} says 21, Maven wins and your build
will behave unexpectedly.

\begin{analogy}{What Maven actually does}
Think of \code{pom.xml} as a shopping list plus a recipe.

The shopping list is your dependencies: ``I need Pi4J version 3.0.2''. Maven
fetches them from a central warehouse and caches them in
\code{\textasciitilde/.m2/}, so the second build is offline and instant.

The recipe is the build: compile these sources with this Java version, run these
tests, put the output there. You do not write build steps; you declare what you
want and Maven knows the order.

The first build downloads a lot and feels slow. That is a one-time cost.
\end{analogy}

\section{git}

\begin{shell}
# macOS -- may already be present via Xcode tools
brew install git

# Debian / Ubuntu / WSL2
sudo apt install -y git
\end{shell}

Tell git who you are. This goes into every commit you ever make, so use a real
address:

\begin{shell}
git config --global user.name  "Your Name"
git config --global user.email "you@example.com"

# Name the default branch 'main' -- matches GitHub and this book
git config --global init.defaultBranch main

# Show the config you just set
git config --global --list
\end{shell}

\begin{trap}{Your commit email becomes public}
If you push to a public GitHub repository, every commit shows the email you set
here, forever, in the git history. It is not editable after the fact without
rewriting history.

If that bothers you, GitHub can give you a private no-reply address: find it
under \emph{Settings $\rightarrow$ Emails} and use that instead. Set it now
rather than discovering it in six months.
\end{trap}

\section{One more useful tool}

\code{rsync} copies files efficiently --- only the parts that changed. You will
use it in Chapter~\ref{ch:deploy} to push builds to the Pi.

\begin{shell}
# macOS: already installed
# Debian / Ubuntu / WSL2
sudo apt install -y rsync

rsync --version | head -1
\end{shell}

\begin{checkpoint}{All four tools}
\begin{shell}
java -version     # 21.x
mvn -v            # 3.9.x, reporting Java 21
git --version     # 2.x
rsync --version   # 3.x
\end{shell}

All four print a version. None says ``command not found''. If Maven reports a
different Java version than \code{java -version} does, fix that now ---
Appendix~\ref{app:trouble0} has the two-line solution.
\end{checkpoint}
